% !TeX root = ../main.tex

\chapter{动态仓位的作用与边界}
\section{风控是否实际触发止损}
\begin{table}[htbp]
\centering
\caption{风控实际仓位与月度风险诊断}
\renewcommand{\arraystretch}{1.18}
\small
\setlength{\tabcolsep}{4.5pt}
\begin{tabularx}{\textwidth}{@{}p{1.55cm}*{4}{C}@{}}
\toprule
路径 & 平均仓位/\% & 零仓位月/个 & 最差单月/\% & 上证下跌月均收益/\% \\
\midrule
P1 & 100.00 & 0 & -30.10 & -3.20 \\
P2 & 100.00 & 0 & -28.82 & -3.28 \\
P3 & 80.40 & 0 & -14.02 & -2.29 \\
P4 & 78.62 & 0 & -13.45 & -2.37 \\
\bottomrule
\end{tabularx}
\end{table}

P3平均仓位80.40\%，P4为78.62\%，两者零仓位月份均为0。按代码机制，若季度止损清仓实际生效，应出现零仓位月份；本样本没有该现象。因此，不能把最大回撤下降解释为25\%季度止损已经成功阻止损失，已观察到的风控作用主要体现为波动目标下的持续缩仓。对止损阈值本身的有效性仍缺少实际触发证据。

\begin{figure}[htbp]
\centering
\includegraphics[width=\linewidth]{fig13}
\caption{动态风控的月度目标仓位及年度平均仓位}
\par\smallskip
\begin{minipage}{0.97\linewidth}\small 注：仓位只使用月初已知历史收益；未加杠杆，现金利息为零。零仓位月数为零。资料来源：CSMAR，作者计算。\end{minipage}
\end{figure}

\section{收益代价与月度尾部风险}
P1最差单月为$-$30.10\%，P3为$-$14.02\%；P2为$-$28.82\%，P4为$-$13.45\%。在上证下跌月份，两组受控收益平均损失也较未控组合收窄。它们支持“缩仓降低本样本月度风险暴露”的描述，仍不能推出月内损失受到严格上限保护。

两条风控路径全期年化均比底层下降约6.5个百分点，说明风险降低不是无成本增加收益。实际部署还会受到现金利息、动态交易费用、持仓估值与可成交性的影响；本次没有引入这些机制，不能把当前缩放净收益当作完整账户回测。

\section{风控收益为什么会下降}
动态仓位在多数高波动月份减少暴露，既会缩小负收益，也会按相同比例缩小随后上涨。波动率依赖过去最多12个月的收益，市场快速反弹时，窗口内仍可能保留此前的高波动记录，仓位不能立刻恢复到1。因此，受控组合与底层组合名单相同，长期财富却可能明显不同；这是仓位机制的结果，不能把全部收益下降说成选股能力退化。

P3与P4的平均仓位分别为80.40\%和78.62\%，但不能把底层CAGR简单乘以平均仓位来预测受控CAGR。每个月的仓位与当月收益组合后再复利，仓位发生在哪些上涨或下跌月份，比单一平均值更重要。实际收益还沿用“底层净收益乘仓位”，基础费用一并缩放，现金利息为零；因此它是当前规则下的收益缩放模拟，而非完整交易账户。

风控改善也不在所有阶段同样显著。例如，2015年7月至2024年12月，P1与P3分段最大回撤均约29.42\%，P2与P4均约34.12\%。这说明完整区间回撤大幅收窄，不能直接推广为每一段下跌都得到等比例保护。附录D分别记录各路径最深回撤所在时段，有助于避免只比较某个历史百分数。
